\documentclass[11pt,a4paper]{article}
\usepackage[utf8]{inputenc}
\usepackage[T1]{fontenc}
\usepackage[brazil]{babel}
\usepackage[a4paper,margin=2.3cm]{geometry}
\usepackage{enumitem}
\usepackage{hyperref}
\usepackage{xcolor}
\setlength{\parindent}{0pt}
\setlength{\parskip}{5pt}
\sloppy
\setlength{\emergencystretch}{3em}
\hypersetup{colorlinks=true,urlcolor=blue,linkcolor=black}
\newcommand{\cabecalho}[2]{\begin{center}{\small \textbf{Sprint 2 -- #1}}\\[1mm]\textit{Inventário de Ativos e Vulnerabilidades}\\Matheus Moreli Durão Silva -- 12621CBS217\\Cibersegurança -- Bacharelado -- Noturno\end{center}\hrule\vspace{4mm}}
\newcommand{\bloco}[1]{\par\vspace{2mm}\textbf{#1}\par}
\newcommand{\codigo}[1]{\texttt{#1}}

\begin{document}\cabecalho{S2\_03 -- O que é CRUD?}{}
\bloco{Significado} CRUD reúne quatro ações sobre registros: \textit{Create} (criar), \textit{Read} (consultar), \textit{Update} (atualizar) e \textit{Delete} (excluir). No sistema, o identificador é a chave do dicionário e não é alterado durante o ciclo.
\bloco{Aplicação} Para ativos: \codigo{cadastrar\_ativo}, \codigo{consultar\_por\_id}, \codigo{atualizar\_ativo} e \codigo{excluir\_ativo}. Para vulnerabilidades: \codigo{cadastrar\_vulnerabilidade}, \codigo{consultar\_vulnerabilidades}, \codigo{atualizar\_vulnerabilidade} e \codigo{excluir\_vulnerabilidade}. A vulnerabilidade é identificada pelo CVE dentro do ativo, portanto não precisa de outro ID.
\bloco{Exemplo de sequência} Cadastrar o ID 10; buscá-lo pelo ID; atualizar o responsável; consultar novamente e verificar que ainda é o ID 10; tentar excluir com \codigo{False} e conferir que ele permanece; excluir com \codigo{True}. Para vulnerabilidade, cadastrar um CVE, mudar seu status e removê-lo somente após confirmação. Para um ID ou CVE ausente, as operações mostram mensagem tratada.
\bloco{Como testar} Os testes automatizados cobrem cadastro, ID duplicado, busca inexistente, atualização com preservação do ID, exclusão confirmada/cancelada, CVE duplicado e atualização/remoção de vulnerabilidade.
\bloco{Checklist atendido} Quatro operações disponíveis, identidade preservada, entradas inválidas rejeitadas, inexistência tratada e ciclo completo testado.
\end{document}
